\documentclass[11pt]{article}
\usepackage[hmargin=1.05in,vmargin=0.8in]{geometry}
\usepackage{amsmath,amssymb}
\usepackage[hidelinks]{hyperref}

\title{A counterexample to Conjecture 00000008785}
\author{Feiyu (Independent researcher)}
\date{}

\begin{document}
\maketitle

\begin{abstract}
The claimed single-step factor $1-2/\sqrt\kappa$ is negative for
the positive definite matrix $\operatorname{diag}(1,3/2,2)$.
For almost every unit start the first relative Ritz error is positive
and the second Lanczos step exists. The claimed factor already fails
at that update.
\end{abstract}

\section*{The single-step claim}
Both \href{https://github.com/The-Last-Math-Competition/The-Last-Math-Competition/blob/main/conjectures/00000008785.md}{versions of the conjecture}
specify the exact single-step rate $1-2/\sqrt\kappa$.
For the largest eigenvalue $\lambda_{\max}$ and largest Ritz value
$\theta_k$ in the kth Krylov subspace, let
$E_k=|\lambda_{\max}-\theta_k|/|\lambda_{\max}|$.
An exact single-step contraction factor for this relative error, or
an exact worst-case factor, implies
$E_{k+1}\le (1-2/\sqrt\kappa)E_k$ whenever $E_k>0$ and the next step exists.

\section*{Counterexample}
Take $A=\operatorname{diag}(1,3/2,2)$. Its largest eigenvalue is $2$,
and its Euclidean condition number is $2$. Let $x$ be a unit starting
vector with first two coordinates $x_0,x_1$ nonzero. The excluded
coordinate sections of the unit sphere have surface measure zero.
In the first Krylov subspace, the Ritz value is
\[
 \theta_1=\frac{x_0^2+(3/2)x_1^2+2x_2^2}{x_0^2+x_1^2+x_2^2}<2 .
\]
Thus $E_1>0$. The vectors $x$ and $Ax$ are linearly independent.
Indeed, if $Ax=cx$, the first coordinate forces $c=1$, while the
second forces $c=3/2$. The second Ritz value is consequently defined
as the largest Rayleigh quotient over nonzero vectors in
$\operatorname{span}\{x,Ax\}$. It exists by compactness of the unit
sphere in that subspace.

Every relative Ritz error is nonnegative. Therefore
$E_2\ge 0>(1-2/\sqrt{2})E_1$, contradicting the proposed factor for
almost every start on the unit sphere.

\section*{Lean verification}
The Lean project uses Mathlib matrices, matrix spectrum, submodules,
and the Euclidean operator norm. It defines the Rayleigh quotient and
the second Krylov subspace directly. The final theorem proves that
the second Krylov vector extends the first subspace and that the
proposed bound fails for every start with $x_0x_1\ne 0$.

Run \texttt{lake build} in \texttt{lean/} with the pinned Lean and
Mathlib v4.33.1. The final theorem depends only on
\texttt{propext}, \texttt{Classical.choice}, and \texttt{Quot.sound}.
\end{document}
